% --- Archon physics formalization source begin ---
% archon:physics
% archon:covers PhyXMiniProblems/problem_phyx_mini_0387.lean
% archon:source-report reports/phyx_mini/problem_phyx_mini_0387.source.json

\chapter{Physics problem phyx\_mini\_0387}
\label{ch:PhyXMiniProblems_problem_phyx_mini_0387}

\paragraph{Problem source.}
\#\# Physical scenario

A U-tube manometer filled with water (density = 1000 kg/m\$\textasciicircum{}3\$) shows a height difference of 25 cm. What is the gauge pressure?

\#\# Question

If the right branch is tilted to make an angle of \$30\textasciicircum{}\{\textbackslash{}circ\}\$ with the horizontal, as shown in figure, what should the length of the column in the tilted tube be relative to the U-tube?

\#\# Answer choices

- A: A: 490cm

- B: B: 50cm

- C: C: 154cm

- D: D: 10.2cm

\#\# Figure caption (auxiliary; use the image as primary evidence)

The image shows two diagrams with a U-shaped tube filled with liquid:

1. **Left Diagram:**
   - A vertical U-shaped tube is partly filled with a blue liquid.
   - The level of the liquid in the left arm is higher than in the right arm.
   - An arrow labeled "h" indicates the difference in liquid height between the two arms, with a dashed horizontal line showing the reference level.

2. **Right Diagram:**
   - A similar U-shaped tube is tilted at a 30-degree angle to the horizontal.
   - The same blue liquid is present, forming a line along the slanted tube.
   - An arrow labeled "L" points along the liquid's surface in the tilted tube.
   - The dashed lines indicate the liquid levels relative to the horizontal, matching the height difference from the left diagram.

Both diagrams illustrate concepts of liquid equilibrium in connected tubes.

\#\# Dataset metadata

Temperature and Heat Transfer; Physical Model Grounding Reasoning, Spatial Relation Reasoning

\paragraph{Recorded answer/context.}
B: B: 50cm

\paragraph{Figure/image path.}
/root/proposal\_for\_physic/science-mango/phyx\_mini\_run/phyx\_data/test\_image/387.png

\paragraph{Formalization target.}
create a compiling Lean file with sorry bodies at `PhyXMiniProblems/problem\_phyx\_mini\_0387.lean`. The Lean declarations must preserve the physical quantities, dimensions or dimensional roles, figure labels, governing-law hypotheses, and final relation expressed by this problem.
Use Mathlib/Physlib names found through LeanExplore where available. If a physics API is missing, introduce faithful local abstractions rather than scalar placeholder aliases.

\begin{theorem}[Physics formalization target]
\label{thm:physics:phyx_mini_0387:target}
\lean{PhyXMiniProblems.ProblemPhyXMini0387.problem_phyx_mini_0387}
\uses{def:physics:phyx-mini-0387:phyxminiproblems-problemphyxmini0387-waterutubemanometersetup, def:physics:phyx-mini-0387:phyxminiproblems-problemphyxmini0387-matchessuppliedmanometerfigure, def:physics:phyx-mini-0387:phyxminiproblems-problemphyxmini0387-matchesproblemreadouts, def:physics:phyx-mini-0387:phyxminiproblems-problemphyxmini0387-hasphysicalmanometerparameters, def:physics:phyx-mini-0387:phyxminiproblems-problemphyxmini0387-satisfieswatermanometerandtiltlaws, def:physics:phyx-mini-0387:phyxminiproblems-problemphyxmini0387-answersinclinedcolumnquestion}
The assigned autoformalize agent should translate this physics problem into Lean declarations in the covered file, with theorem and lemma proof bodies written as `by sorry`.
\end{theorem}
\begin{proof}
This is an autoformalization task, not a proof task. Produce statements that can later be proved by the physics prover without weakening the physical model.
\end{proof}

% --- Archon declaration topology begin ---
\section{Declaration topology}
\begin{definition}[mass Density Dimension]
  \label{def:physics:phyx-mini-0387:phyxminiproblems-problemphyxmini0387-massdensitydimension}
  \lean{PhyXMiniProblems.ProblemPhyXMini0387.massDensityDimension}
  The physical dimension of mass density, M L⁻³
\end{definition}
\begin{proof}
  This is the declared model definition of mass Density Dimension.
\end{proof}
\begin{definition}[acceleration Dimension]
  \label{def:physics:phyx-mini-0387:phyxminiproblems-problemphyxmini0387-accelerationdimension}
  \lean{PhyXMiniProblems.ProblemPhyXMini0387.accelerationDimension}
  The physical dimension of acceleration, L T⁻²
\end{definition}
\begin{proof}
  This is the declared model definition of acceleration Dimension.
\end{proof}
\begin{definition}[Length Quantity]
  \label{def:physics:phyx-mini-0387:phyxminiproblems-problemphyxmini0387-lengthquantity}
  \lean{PhyXMiniProblems.ProblemPhyXMini0387.LengthQuantity}
  A nonnegative physical length, independent of the selected readout unit
\end{definition}
\begin{proof}
  This is the declared model definition of Length Quantity.
\end{proof}
\begin{definition}[Mass Density Quantity]
  \label{def:physics:phyx-mini-0387:phyxminiproblems-problemphyxmini0387-massdensityquantity}
  \lean{PhyXMiniProblems.ProblemPhyXMini0387.MassDensityQuantity}
  \uses{def:physics:phyx-mini-0387:phyxminiproblems-problemphyxmini0387-massdensitydimension}
  A nonnegative physical mass density
\end{definition}
\begin{proof}
  This is the declared model definition of Mass Density Quantity.
\end{proof}
\begin{definition}[Acceleration Quantity]
  \label{def:physics:phyx-mini-0387:phyxminiproblems-problemphyxmini0387-accelerationquantity}
  \lean{PhyXMiniProblems.ProblemPhyXMini0387.AccelerationQuantity}
  \uses{def:physics:phyx-mini-0387:phyxminiproblems-problemphyxmini0387-accelerationdimension}
  A nonnegative physical acceleration magnitude
\end{definition}
\begin{proof}
  This is the declared model definition of Acceleration Quantity.
\end{proof}
\begin{definition}[Pressure Quantity]
  \label{def:physics:phyx-mini-0387:phyxminiproblems-problemphyxmini0387-pressurequantity}
  \lean{PhyXMiniProblems.ProblemPhyXMini0387.PressureQuantity}
  Gauge or absolute pressure, using Physlib's pressure dimension
\end{definition}
\begin{proof}
  This is the declared model definition of Pressure Quantity.
\end{proof}
\begin{definition}[length Readout]
  \label{def:physics:phyx-mini-0387:phyxminiproblems-problemphyxmini0387-lengthreadout}
  \lean{PhyXMiniProblems.ProblemPhyXMini0387.lengthReadout}
  \uses{def:physics:phyx-mini-0387:phyxminiproblems-problemphyxmini0387-lengthquantity}
  Read a physical length in a selected length unit
\end{definition}
\begin{proof}
  This is the declared model definition of length Readout.
\end{proof}
\begin{definition}[length In Meters]
  \label{def:physics:phyx-mini-0387:phyxminiproblems-problemphyxmini0387-lengthinmeters}
  \lean{PhyXMiniProblems.ProblemPhyXMini0387.lengthInMeters}
  \uses{def:physics:phyx-mini-0387:phyxminiproblems-problemphyxmini0387-lengthquantity, def:physics:phyx-mini-0387:phyxminiproblems-problemphyxmini0387-lengthreadout}
  Metre readout of a physical length
\end{definition}
\begin{proof}
  This is the declared model definition of length In Meters.
\end{proof}
\begin{definition}[length In Centimeters]
  \label{def:physics:phyx-mini-0387:phyxminiproblems-problemphyxmini0387-lengthincentimeters}
  \lean{PhyXMiniProblems.ProblemPhyXMini0387.lengthInCentimeters}
  \uses{def:physics:phyx-mini-0387:phyxminiproblems-problemphyxmini0387-lengthquantity, def:physics:phyx-mini-0387:phyxminiproblems-problemphyxmini0387-lengthreadout}
  Centimetre readout of a physical length
\end{definition}
\begin{proof}
  This is the declared model definition of length In Centimeters.
\end{proof}
\begin{definition}[density In Kilograms Per Cubic Meter]
  \label{def:physics:phyx-mini-0387:phyxminiproblems-problemphyxmini0387-densityinkilogramspercubicmeter}
  \lean{PhyXMiniProblems.ProblemPhyXMini0387.densityInKilogramsPerCubicMeter}
  \uses{def:physics:phyx-mini-0387:phyxminiproblems-problemphyxmini0387-massdensityquantity}
  Kilogram-per-cubic-metre readout of a physical mass density
\end{definition}
\begin{proof}
  This is the declared model definition of density In Kilograms Per Cubic Meter.
\end{proof}
\begin{definition}[acceleration In Meters Per Second Squared]
  \label{def:physics:phyx-mini-0387:phyxminiproblems-problemphyxmini0387-accelerationinmeterspersecondsquared}
  \lean{PhyXMiniProblems.ProblemPhyXMini0387.accelerationInMetersPerSecondSquared}
  \uses{def:physics:phyx-mini-0387:phyxminiproblems-problemphyxmini0387-accelerationquantity}
  Metre-per-second-squared readout of an acceleration magnitude
\end{definition}
\begin{proof}
  This is the declared model definition of acceleration In Meters Per Second Squared.
\end{proof}
\begin{definition}[pressure In Pascals]
  \label{def:physics:phyx-mini-0387:phyxminiproblems-problemphyxmini0387-pressureinpascals}
  \lean{PhyXMiniProblems.ProblemPhyXMini0387.pressureInPascals}
  \uses{def:physics:phyx-mini-0387:phyxminiproblems-problemphyxmini0387-pressurequantity}
  Pascal readout of a pressure
\end{definition}
\begin{proof}
  This is the declared model definition of pressure In Pascals.
\end{proof}
\begin{definition}[degrees]
  \label{def:physics:phyx-mini-0387:phyxminiproblems-problemphyxmini0387-degrees}
  \lean{PhyXMiniProblems.ProblemPhyXMini0387.degrees}
  Convert a numerical degree readout to Mathlib's physical angle type
\end{definition}
\begin{proof}
  This is the declared model definition of degrees.
\end{proof}
\begin{definition}[Manometer Liquid]
  \label{def:physics:phyx-mini-0387:phyxminiproblems-problemphyxmini0387-manometerliquid}
  \lean{PhyXMiniProblems.ProblemPhyXMini0387.ManometerLiquid}
  The liquid filling the connected lower part of the manometer
\end{definition}
\begin{proof}
  This is the declared model definition of Manometer Liquid.
\end{proof}
\begin{definition}[Manometer Arm]
  \label{def:physics:phyx-mini-0387:phyxminiproblems-problemphyxmini0387-manometerarm}
  \lean{PhyXMiniProblems.ProblemPhyXMini0387.ManometerArm}
  The two arms, with the right arm being the one tilted in the figure
\end{definition}
\begin{proof}
  This is the declared model definition of Manometer Arm.
\end{proof}
\begin{definition}[Figure Panel]
  \label{def:physics:phyx-mini-0387:phyxminiproblems-problemphyxmini0387-figurepanel}
  \lean{PhyXMiniProblems.ProblemPhyXMini0387.FigurePanel}
  The two side-by-side configurations drawn in the primary image
\end{definition}
\begin{proof}
  This is the declared model definition of Figure Panel.
\end{proof}
\begin{definition}[Figure Length Label]
  \label{def:physics:phyx-mini-0387:phyxminiproblems-problemphyxmini0387-figurelengthlabel}
  \lean{PhyXMiniProblems.ProblemPhyXMini0387.FigureLengthLabel}
  The two physical length labels printed in the primary image
\end{definition}
\begin{proof}
  This is the declared model definition of Figure Length Label.
\end{proof}
\begin{definition}[Angle Reference]
  \label{def:physics:phyx-mini-0387:phyxminiproblems-problemphyxmini0387-anglereference}
  \lean{PhyXMiniProblems.ProblemPhyXMini0387.AngleReference}
  The baseline from which the displayed 30° inclination is measured
\end{definition}
\begin{proof}
  This is the declared model definition of Angle Reference.
\end{proof}
\begin{definition}[Supplied Manometer Figure]
  \label{def:physics:phyx-mini-0387:phyxminiproblems-problemphyxmini0387-suppliedmanometerfigure}
  \lean{PhyXMiniProblems.ProblemPhyXMini0387.SuppliedManometerFigure}
  \uses{def:physics:phyx-mini-0387:phyxminiproblems-problemphyxmini0387-lengthquantity, def:physics:phyx-mini-0387:phyxminiproblems-problemphyxmini0387-manometerliquid, def:physics:phyx-mini-0387:phyxminiproblems-problemphyxmini0387-manometerarm, def:physics:phyx-mini-0387:phyxminiproblems-problemphyxmini0387-figurepanel, def:physics:phyx-mini-0387:phyxminiproblems-problemphyxmini0387-figurelengthlabel, def:physics:phyx-mini-0387:phyxminiproblems-problemphyxmini0387-anglereference}
  A structured transcription of image 387.png. The labelled lengths are independent physical quantities. In particular, L is not defined to be 50 cm; its relation to h is supplied only by the geometric governing law below
\end{definition}
\begin{proof}
  This is the declared model definition of Supplied Manometer Figure.
\end{proof}
\begin{definition}[Water UTube Manometer Setup]
  \label{def:physics:phyx-mini-0387:phyxminiproblems-problemphyxmini0387-waterutubemanometersetup}
  \lean{PhyXMiniProblems.ProblemPhyXMini0387.WaterUTubeManometerSetup}
  \uses{def:physics:phyx-mini-0387:phyxminiproblems-problemphyxmini0387-massdensityquantity, def:physics:phyx-mini-0387:phyxminiproblems-problemphyxmini0387-accelerationquantity, def:physics:phyx-mini-0387:phyxminiproblems-problemphyxmini0387-pressurequantity, def:physics:phyx-mini-0387:phyxminiproblems-problemphyxmini0387-manometerarm, def:physics:phyx-mini-0387:phyxminiproblems-problemphyxmini0387-suppliedmanometerfigure}
  The physical manometer, its boundary pressures, and the quantities in the figure
\end{definition}
\begin{proof}
  This is the declared model definition of Water UTube Manometer Setup.
\end{proof}
\begin{definition}[height Difference H]
  \label{def:physics:phyx-mini-0387:phyxminiproblems-problemphyxmini0387-heightdifferenceh}
  \lean{PhyXMiniProblems.ProblemPhyXMini0387.heightDifferenceH}
  \uses{def:physics:phyx-mini-0387:phyxminiproblems-problemphyxmini0387-lengthquantity, def:physics:phyx-mini-0387:phyxminiproblems-problemphyxmini0387-waterutubemanometersetup}
  The physical height difference denoted by h in the upright panel
\end{definition}
\begin{proof}
  This is the declared model definition of height Difference H.
\end{proof}
\begin{definition}[tilted Column Length L]
  \label{def:physics:phyx-mini-0387:phyxminiproblems-problemphyxmini0387-tiltedcolumnlengthl}
  \lean{PhyXMiniProblems.ProblemPhyXMini0387.tiltedColumnLengthL}
  \uses{def:physics:phyx-mini-0387:phyxminiproblems-problemphyxmini0387-lengthquantity, def:physics:phyx-mini-0387:phyxminiproblems-problemphyxmini0387-waterutubemanometersetup}
  The physical inclined liquid-column length denoted by L
\end{definition}
\begin{proof}
  This is the declared model definition of tilted Column Length L.
\end{proof}
\begin{definition}[Matches Supplied Manometer Figure]
  \label{def:physics:phyx-mini-0387:phyxminiproblems-problemphyxmini0387-matchessuppliedmanometerfigure}
  \lean{PhyXMiniProblems.ProblemPhyXMini0387.MatchesSuppliedManometerFigure}
  \uses{def:physics:phyx-mini-0387:phyxminiproblems-problemphyxmini0387-waterutubemanometersetup}
  Qualitative information read directly from the supplied bitmap. The lower left surface and higher right surface fix the pressure-sign convention, but no numerical value for L occurs here
\end{definition}
\begin{proof}
  This is the declared model definition of Matches Supplied Manometer Figure.
\end{proof}
\begin{definition}[Matches Problem Readouts]
  \label{def:physics:phyx-mini-0387:phyxminiproblems-problemphyxmini0387-matchesproblemreadouts}
  \lean{PhyXMiniProblems.ProblemPhyXMini0387.MatchesProblemReadouts}
  \uses{def:physics:phyx-mini-0387:phyxminiproblems-problemphyxmini0387-lengthincentimeters, def:physics:phyx-mini-0387:phyxminiproblems-problemphyxmini0387-densityinkilogramspercubicmeter, def:physics:phyx-mini-0387:phyxminiproblems-problemphyxmini0387-degrees, def:physics:phyx-mini-0387:phyxminiproblems-problemphyxmini0387-waterutubemanometersetup, def:physics:phyx-mini-0387:phyxminiproblems-problemphyxmini0387-heightdifferenceh}
  Numerical data stated in the problem and printed in the figure: water has the given density, the upright head is 25 cm, and the right branch is inclined at 30°. The requested length L is deliberately absent
\end{definition}
\begin{proof}
  This is the declared model definition of Matches Problem Readouts.
\end{proof}
\begin{definition}[Has Physical Manometer Parameters]
  \label{def:physics:phyx-mini-0387:phyxminiproblems-problemphyxmini0387-hasphysicalmanometerparameters}
  \lean{PhyXMiniProblems.ProblemPhyXMini0387.HasPhysicalManometerParameters}
  \uses{def:physics:phyx-mini-0387:phyxminiproblems-problemphyxmini0387-lengthinmeters, def:physics:phyx-mini-0387:phyxminiproblems-problemphyxmini0387-densityinkilogramspercubicmeter, def:physics:phyx-mini-0387:phyxminiproblems-problemphyxmini0387-accelerationinmeterspersecondsquared, def:physics:phyx-mini-0387:phyxminiproblems-problemphyxmini0387-pressureinpascals, def:physics:phyx-mini-0387:phyxminiproblems-problemphyxmini0387-waterutubemanometersetup, def:physics:phyx-mini-0387:phyxminiproblems-problemphyxmini0387-heightdifferenceh, def:physics:phyx-mini-0387:phyxminiproblems-problemphyxmini0387-tiltedcolumnlengthl}
  Positivity and nondegeneracy conditions for the physical apparatus
\end{definition}
\begin{proof}
  This is the declared model definition of Has Physical Manometer Parameters.
\end{proof}
\begin{definition}[Satisfies Water Manometer And Tilt Laws]
  \label{def:physics:phyx-mini-0387:phyxminiproblems-problemphyxmini0387-satisfieswatermanometerandtiltlaws}
  \lean{PhyXMiniProblems.ProblemPhyXMini0387.SatisfiesWaterManometerAndTiltLaws}
  \uses{def:physics:phyx-mini-0387:phyxminiproblems-problemphyxmini0387-lengthinmeters, def:physics:phyx-mini-0387:phyxminiproblems-problemphyxmini0387-densityinkilogramspercubicmeter, def:physics:phyx-mini-0387:phyxminiproblems-problemphyxmini0387-accelerationinmeterspersecondsquared, def:physics:phyx-mini-0387:phyxminiproblems-problemphyxmini0387-pressureinpascals, def:physics:phyx-mini-0387:phyxminiproblems-problemphyxmini0387-waterutubemanometersetup, def:physics:phyx-mini-0387:phyxminiproblems-problemphyxmini0387-heightdifferenceh, def:physics:phyx-mini-0387:phyxminiproblems-problemphyxmini0387-tiltedcolumnlengthl}
  Governing pressure and geometry relations. The right opening is the atmospheric reference. Gauge pressure is applied pressure minus atmospheric pressure. Hydrostatic balance gives Δp = ρ g h in coherent SI readouts. The vertical projection of the tilted length is L sin θ = h. These are physical laws, not solved numerical values for L
\end{definition}
\begin{proof}
  This is the declared model definition of Satisfies Water Manometer And Tilt Laws.
\end{proof}
\begin{theorem}[gauge pressure eq density times gravity times height]
  \label{thm:physics:phyx-mini-0387:phyxminiproblems-problemphyxmini0387-gauge-pressure-eq-density-times-gravity-times-height}
  \lean{PhyXMiniProblems.ProblemPhyXMini0387.gauge_pressure_eq_density_times_gravity_times_height}
  \uses{def:physics:phyx-mini-0387:phyxminiproblems-problemphyxmini0387-lengthinmeters, def:physics:phyx-mini-0387:phyxminiproblems-problemphyxmini0387-densityinkilogramspercubicmeter, def:physics:phyx-mini-0387:phyxminiproblems-problemphyxmini0387-accelerationinmeterspersecondsquared, def:physics:phyx-mini-0387:phyxminiproblems-problemphyxmini0387-pressureinpascals, def:physics:phyx-mini-0387:phyxminiproblems-problemphyxmini0387-waterutubemanometersetup, def:physics:phyx-mini-0387:phyxminiproblems-problemphyxmini0387-heightdifferenceh, def:physics:phyx-mini-0387:phyxminiproblems-problemphyxmini0387-satisfieswatermanometerandtiltlaws}
  The hydrostatic gauge-pressure relation for the stated water head
\end{theorem}
\begin{proof}
  The conclusion follows from the governing relations and figure readouts named in the statement of gauge pressure eq density times gravity times height.
\end{proof}
\begin{theorem}[tilted column length is twice height]
  \label{thm:physics:phyx-mini-0387:phyxminiproblems-problemphyxmini0387-tilted-column-length-is-twice-height}
  \lean{PhyXMiniProblems.ProblemPhyXMini0387.tilted_column_length_is_twice_height}
  \uses{def:physics:phyx-mini-0387:phyxminiproblems-problemphyxmini0387-lengthincentimeters, def:physics:phyx-mini-0387:phyxminiproblems-problemphyxmini0387-waterutubemanometersetup, def:physics:phyx-mini-0387:phyxminiproblems-problemphyxmini0387-heightdifferenceh, def:physics:phyx-mini-0387:phyxminiproblems-problemphyxmini0387-tiltedcolumnlengthl, def:physics:phyx-mini-0387:phyxminiproblems-problemphyxmini0387-matchessuppliedmanometerfigure, def:physics:phyx-mini-0387:phyxminiproblems-problemphyxmini0387-matchesproblemreadouts, def:physics:phyx-mini-0387:phyxminiproblems-problemphyxmini0387-hasphysicalmanometerparameters, def:physics:phyx-mini-0387:phyxminiproblems-problemphyxmini0387-satisfieswatermanometerandtiltlaws}
  At 30°, the inclined column is twice the original vertical head
\end{theorem}
\begin{proof}
  The conclusion follows from the governing relations and figure readouts named in the statement of tilted column length is twice height.
\end{proof}
\begin{definition}[Answer Choice]
  \label{def:physics:phyx-mini-0387:phyxminiproblems-problemphyxmini0387-answerchoice}
  \lean{PhyXMiniProblems.ProblemPhyXMini0387.AnswerChoice}
  The four length choices printed in the source problem
\end{definition}
\begin{proof}
  This is the declared model definition of Answer Choice.
\end{proof}
\begin{definition}[length In Centimeters]
  \label{def:physics:phyx-mini-0387:phyxminiproblems-problemphyxmini0387-answerchoice-lengthincentimeters}
  \lean{PhyXMiniProblems.ProblemPhyXMini0387.AnswerChoice.lengthInCentimeters}
  \uses{def:physics:phyx-mini-0387:phyxminiproblems-problemphyxmini0387-lengthincentimeters, def:physics:phyx-mini-0387:phyxminiproblems-problemphyxmini0387-answerchoice}
  Centimetre value displayed beside an answer-choice label
\end{definition}
\begin{proof}
  This is the declared model definition of length In Centimeters.
\end{proof}
\begin{definition}[Answers Inclined Column Question]
  \label{def:physics:phyx-mini-0387:phyxminiproblems-problemphyxmini0387-answersinclinedcolumnquestion}
  \lean{PhyXMiniProblems.ProblemPhyXMini0387.AnswersInclinedColumnQuestion}
  \uses{def:physics:phyx-mini-0387:phyxminiproblems-problemphyxmini0387-lengthincentimeters, def:physics:phyx-mini-0387:phyxminiproblems-problemphyxmini0387-waterutubemanometersetup, def:physics:phyx-mini-0387:phyxminiproblems-problemphyxmini0387-tiltedcolumnlengthl, def:physics:phyx-mini-0387:phyxminiproblems-problemphyxmini0387-answerchoice}
  A choice matches the independently modeled inclined physical length
\end{definition}
\begin{proof}
  This is the declared model definition of Answers Inclined Column Question.
\end{proof}
\begin{theorem}[tilted column length is fifty centimeters]
  \label{thm:physics:phyx-mini-0387:phyxminiproblems-problemphyxmini0387-tilted-column-length-is-fifty-centimeters}
  \lean{PhyXMiniProblems.ProblemPhyXMini0387.tilted_column_length_is_fifty_centimeters}
  \uses{def:physics:phyx-mini-0387:phyxminiproblems-problemphyxmini0387-lengthincentimeters, def:physics:phyx-mini-0387:phyxminiproblems-problemphyxmini0387-waterutubemanometersetup, def:physics:phyx-mini-0387:phyxminiproblems-problemphyxmini0387-tiltedcolumnlengthl, def:physics:phyx-mini-0387:phyxminiproblems-problemphyxmini0387-matchessuppliedmanometerfigure, def:physics:phyx-mini-0387:phyxminiproblems-problemphyxmini0387-matchesproblemreadouts, def:physics:phyx-mini-0387:phyxminiproblems-problemphyxmini0387-hasphysicalmanometerparameters, def:physics:phyx-mini-0387:phyxminiproblems-problemphyxmini0387-satisfieswatermanometerandtiltlaws}
  The requested relation in the column marked L is 50 cm
\end{theorem}
\begin{proof}
  The conclusion follows from the governing relations and figure readouts named in the statement of tilted column length is fifty centimeters.
\end{proof}
% --- Archon declaration topology end ---
% --- Archon physics formalization source end ---
